\section{Supplementary: Feature Importance Analysis}

\label{app:features}
% ══════════════════════════════════════════════════════════════════════════════

We analyze feature importance using integrated gradients
\citep{chen2020deep} applied to the feature encoder. Across all 847
species, the most influential variables are:

\begin{table}[H]
\centering
\caption{Top 10 features by mean integrated gradient magnitude across all
species.}
\label{tab:features}
\begin{tabular}{clc}
\toprule
\textbf{Rank} & \textbf{Feature} & \textbf{Importance} \\
\midrule
1  & Annual mean temperature (BIO1)      & 0.187 \\
2  & Annual precipitation (BIO12)        & 0.156 \\
3  & Elevation                           & 0.134 \\
4  & Tree cover fraction                 & 0.121 \\
5  & Temperature seasonality (BIO4)      & 0.098 \\
6  & NDVI mean                           & 0.089 \\
7  & Precipitation seasonality (BIO15)   & 0.076 \\
8  & Distance to water                   & 0.062 \\
9  & Slope                               & 0.048 \\
10 & Temperature annual range (BIO7)     & 0.041 \\
\bottomrule
\end{tabular}
\end{table}

Feature importance varies substantially across taxa: for amphibians,
precipitation variables dominate; for birds, temperature and NDVI are
most influential; for mammals, tree cover and elevation are primary
drivers.

